\selectlanguage{ngerman}%
\frenchspacing%
\begin{ledgroupsized}[r]{120mm}%
\footnotesize%
\pstart%
\noindent%
\textbf{Überlieferung:}%
\pend%
\end{ledgroupsized}%
\begin{ledgroupsized}[r]{114mm}%
\footnotesize%
\pstart%
\parindent -6mm%
\makebox[6mm][l]{\textit{L}}%
Aufzeichnung:
LH XXXVII 5 Bl.\;121.
Ein Blatt~8\textsuperscript{o},
Ränder beschnitten;
Papiererhaltungsmaßnahmen.
Eine volle Seite auf Bl.\;\href{\Scan{LH_37_05_121r}}{121\,r\textsuperscript{o}}.
Bl.\;\href{\Scan{LH_37_05_121v}}{121\,v\textsuperscript{o}} überliefert lediglich ein Satzfragment von Leibnizens Hand (\textit{Si je n'esperois que}).
\pend%
\end{ledgroupsized}%
\vspace{5mm}%
\begin{ledgroup}%
\footnotesize%
\pstart%
\noindent%
\textbf{Datierungsgründe:}
Die Datierung der Aufzeichnung N.\;\ref{60337} beruht auf Eigenschaften des Textträgers sowie auf inhaltlichen Betrachtungen.
Siehe die ausführliche Begründung auf S.\;\pageref{60337_dat}.%
\pend%
\end{ledgroup}%
\selectlanguage{latin}%
\frenchspacing%
\vspace{8mm}%
\normalsize%
\pstart%
\noindent%
\lbrack%
\href{\Scan{LH_37_05_121r}}%
{121~r\textsuperscript{o}}%
\rbrack\
Centrum tensionis%
\indexS{centrum tensionis}%
\indexS{tensio inaequalis}
in linea,%
\indexS{linea inaequaliter tensa}
si tendatur inaequaliter
ob inaequalitatem materiae,%
\indexS{materia inaequalis}%
\indexS{inaequalitas materiae}
est illud punctum %
cujus tensio ducta in longitudinem,%
\indexS{tensio ducta in longitudinem}
ita ut tota linea uniformiter puncti tensione%
\indexS{tensio puncti}
ponatur%
\indexS{linea uniformiter tensa}
\edtext{%
tensa,
det productum aequale%
}{%
\lemma{tensa,}%
\Bfootnote{%
\textit{(1)}~aeque %
\textit{(2)}~det productum aequale%
~\textit{L}%
}}
facto ex praesentibus tensionibus%
\indexS{tensio praesens}
in lineam ordinatim%
\indexS{tensio ducta in lineam}
\edtext{%
ductis.
Atque ita est centrum%
\indexS{centrum tensionis}
tensionis mediae.%
\indexS{tensio media}
Tensio media autem illa est,
cui si attribuatur longitudo lineae,
idem prodit ac
si cuilibet puncto sua tensio attribuatur:%
\indexS{tensio puncti}
est ergo media arithmetica inter tensiones.%
\indexS{media arithmetica tensionum}
Similiter in centro gravitatis.%
\indexS{centrum gravitatis}%
}{%
\lemma{ductis.}%
\Bfootnote{%
\textit{(1)}~Ubi centrum gravitatis est, cui si attribuatur pondus %
\textit{(2)}~Atque ita %
\lbrack...\rbrack\ centro gravitatis.%
~\textit{L}%
}}
\pend%
\pstart%
\edtext{%
Sit
\edtext{%
recta horizontalis \textit{AGBCP}%
}{%
\lemma{recta horizontalis \textit{AGBCP}}%
\Cfootnote{%
Vgl. das Diagramm \lbrack\textit{Fig.\:1}\rbrack.%
}}
in qua recta gravis \textit{AGB},%
\indexS{recta gravis}
cujus centrum gravitatis \textit{G},%
\indexS{centrum gravitatis}%
}{%
\lemma{Sit}%
\Bfootnote{%
\textit{(1)}~centrum aequ %
\textit{(2)}~recta horizontalis \textit{AGBCP} %
\textit{(a)}~. Sit centru %
\textit{(b)}~in qua %
\lbrack...\rbrack\ centrum gravitatis \textit{G},%
~\textit{L}%
}}
contrapondium \textit{P},%
\indexS{contrapondium}
ita ut maneat aequilibrium.%
\indexS{aequilibrium rectae gravis}
Dico centrum gravitatis \textit{G}%
\indexS{centrum gravitatis}
habere gravitationem mediam,%
\indexS{gravitatio media}
ita
\edtext{%
ut idem momentum%
\indexS{momentum e gravitationibus}
seu factum ex gravitationibus
ad gravis partes ordinatim applicatis
prodeat,%
}{%
\lemma{ut}%
\Bfootnote{%
\textit{(1)}~idem pondus \textit{P} %
\textit{(2)}~eade %
\textit{(3)}~idem momentum %
\lbrack...\rbrack\ gravitationibus ad %
\textit{(a)}~graves ordi %
\textit{(b)}~gravis partes ordinatim applicatis prodeat,%
~\textit{L}%
}}
sive cuilibet puncto rectae \textit{AB}
suam tribuas gravitationem,%
\indexS{gravitatio puncti}
sive omnibus punctis rectae \textit{AB}
tribuas gravitationem mediam%
\indexS{gravitatio media}
quae est puncti \textit{G}.
Quod fit
si \textit{AB} colloces in situ \textit{(A)(B)},
vel si libere suspendas ex \textit{G}.
\pend%
\pstart%
\edlabel{LH_37_05_121_121r_bfs-1}%
Uti
\edtext{%
consideratio centri gravitatis%
\indexS{consideratio centri gravitatis}%
\indexS{centrum gravitatis}
\edtext{%
prodest
ad motum perpetuum refutandum,%
\indexS{motus perpetuus refutandus}%
}{%
\lemma{prodest}%
\Bfootnote{%
\textit{(1)}~contra %
\textit{(2)}~ad motum perpetuum refutandum,%
~\textit{L}%
}}
quia ut motus vi gravitatis%
\indexS{vis gravitatis}
fiat,%
\indexS{motus factus vi gravitatis}
oportet centrum gravitatis%
\indexS{centrum gravitatis}
totale descendere;%
\indexS{descensus centri gravitatis}%
}{%
\lemma{consideratio \lbrack...\rbrack\ descendere}%
\Cfootnote{%
Siehe zu diesem Grundsatz
\textit{LSB} VIII,\;3 N.\;39, S.\;387.16--388.6;\cite{01488}
N.\;58\textsubscript{11}, S.\;645.10--12\cite{01401}
(samt Quellenangaben in den dazugehörigen Erläuterungen).%
}}
ita
\edtext{%
e\lbrack a\rbrack dem%
}{%
\lemma{eodem}%
\Bfootnote{%
\textit{L\;ändert Hrsg.}%
}}
proderit etiam consideratio centri tensionis,%
\indexS{consideratio centri tensionis}%
\indexS{centrum tensionis}
nam non potest fieri operatio%
\indexS{operatio facta vi tensionis}
vi tensionis,%
\indexS{vis tensionis}
quin tensio media%
\indexS{tensio media}
minuatur
\edtext{%
seu laxetur.%
\edlabel{LH_37_05_121_121r_bfs-2}%
}{%
\lemma{seu laxetur}%
\Bfootnote{%
\textit{erg.\;L}%
}}
\ledpb%
Quod si conjungantur haec duo,
gravitas%
\indexS{gravitas conjuncta tensioni}
et tensio,%
\indexS{tensio conjuncta gravitati}
non potest fieri motus nisi
\edtext{%
plus descendatur%
}{%
\lemma{plus}%
\Bfootnote{%
\textit{(1)}~gravit\textlangle a\textrangle s %
\textit{(2)}~descendatur%
~\textit{L}%
}}
quam laxetur;
comparando scilicet ambo inter se.
Idem erit si accedat directio magnetica.%
\indexS{directio magnetica}
Postremo impetus ipse%
\indexS{impetus conjunctus tensioni}%
\indexS{impetus conjunctus gravitati}
\edtext{%
his conjungi potest,%
}{%
\lemma{his}%
\Bfootnote{%
\textit{(1)}~conjungitur %
\textit{(2)}~conjungi potest,%
~\textit{L}%
}}
pro quo substituamus tensionem%
\indexS{tensio substituta impetui}
vel elevationem%
\indexS{elevatio substituta impetui}
quam producere potest.
\pend%
\begin{edfigure}[2.50em][0.00em]
\includegraphics%
{gesamttex/edit_VIII-4/images/LH_37_05_121_d_121r.pdf}
\edcaption[0.0em]{%
\lbrack\textit{Fig.\:1}\rbrack%
\edtext{}{%
\lemma{\lbrack\textit{Fig.\:1}\rbrack}%
\Cfootnote{%
In \textit{L} (Bl.\;\href{\Scan{LH_37_05_121r}}{121\,r\textsuperscript{o}}) ist eine gestr., hier nicht wiedergegebene Vorstufe des Diagramms  überliefert.%
}}%
}%
\end{edfigure}%
\vspace{3.0em}%
\pstart%
\noindent%
\centering%
*     *     *
\pend%
\vspace{2.5em}%
\selectlanguage{german}%
\frenchspacing%
\footnotesize%
\pstart%
\noindent%
\textls{Ausführliche Datierungsgründe:}%
\pagelabel{60337_dat}
Der Träger der titellosen und undatierten Aufzeichnung N.\;\ref{60337} überliefert auf der Rücksiete (Bl.\;\href{\Scan{LH_37_05_121v}}{121\,v\textsuperscript{o}}) ein Satzfragment von Leibnizens Hand (\textit{Si je n'esperois que}), welches eine gewisse Ähnlichkeit mit dem Anfang von Leibnizens Brief an Guillaume François de L'Hospital vom 20./30. Juli 1696 aufweist:
\textit{Je serois fort en peine de Vostre santé, si je n'esperois, Monsieur, que vostre silence doit venir de ce que ...} 
(\textit{LSB} III,\;7 N.\;11, S.\;39.20-21).\cite{01483}
Möglicherweise hatte Leibniz begonnen, seinen Brief an L'Hospital auf dem Blatt, auf dem er N.\;\ref{60337} bereits verfasst hatte oder später verfassen würde, zu entwerfen, gab aber gleich seine Absicht auf.
Diese nicht weiter begründbare Annahme ermöglicht die Vermutung, dass die vorliegende Aufzeichnung kurz vor oder nach dem 20./30. Juli 1696 entstand.
Eine frühere oder spätere Entstehungszeit ist nicht auzuschließen.
\pend%
\pstart%
Schließlich ist zu bemerken, dass die im Text angedeutete Analogie zwischen \textit{centrum tensionis} und Schwerpunkt (\refpassage{LH_37_05_121_121r_bfs-1}{LH_37_05_121_121r_bfs-2}) Erkenntnisse der Statik voraussetzt, die Leibniz bereits in seinen frühen Hannoveraner\indexO{Hannover} Jahren besaß und wohl während seines Pariser Aufenthalts erworben hatte;
siehe \textit{LSB} VIII,\;3 N.\;39, S.\;387.16--388.6 (März 1677)\cite{01488}
und \textit{LSB} VIII,\;3 N.\;58\textsubscript{11}, S.\;645.10--12\cite{01401} (Januar-Februar 1678)
samt Quellenangaben in den dazugehörigen Erläuterungen.%
\pend%
\selectlanguage{latin}%
\frenchspacing%
\newpage%
